% SGA 3，第 VII_B 篇，英文重建稿——第 1 轮。
% 范围：权威底本分篇第 97--99 页，末尾参考文献全文，自 [CA] 至 [We95]。
% 合并阅读本页码： 610--612.
% 再版印刷页码： 599--601.
% 本范围内编者注： 170.
% 权威底本： Polo--Gille Exp7B-13oct24.pdf,
% SHA-256 D303690E0195654BBE7C3DAA5D89EA23E24CD709A3575217234AC19AE8836694.
% 曾参考 Jacob Reinhold 的 CC BY 4.0 Markdown 作为比对谱系。
% 书目题名、作者姓名拼写、期刊缩写和出版信息均遵循 Polo--Gille 权威底本。

\chapter*{参考文献}
\phantomsection
\label{sga3:VIIB:bibliography}
\addcontentsline{toc}{chapter}{参考文献}

\begin{description}
  \item[\textup{[CA]}] P.~Gabriel，
  \emph{阿贝尔范畴〔Des cat\'egories ab\'eliennes〕}，
  Bull.\ Soc.\ Math.\ France \textbf{90} (1962)，323--448。
  \footnote{\label{sga3:VIIB:nde:170}编者在原文已有的此条参考文献之后增加了以下各条参考文献。}

  \item[\textup{[Ab80]}] E.~Abe，
  \emph{Hopf 代数〔Hopf algebras〕}，Cambridge Univ.\ Press，1980。

  \item[\textup{[Br00]}] C.~Breuil，
  \emph{p-可除群、有限群与滤过模〔Groupes \(p\)-divisibles, groupes finis et modules filtr\'es〕}，Ann.\ of Math.\ \textbf{152} (2000)，第~2~期，489--549。

  \item[\textup{[BAlg]}] N.~Bourbaki，
  \emph{代数〔Alg\`ebre〕}，第 I--III 及 IV--VII 章，Hermann，1970，及 Masson，1981。

  \item[\textup{[BAC]}] N.~Bourbaki，
  \emph{交换代数〔Alg\`ebre commutative〕}，第 I--IV 章，Masson，1985。

  \item[\textup{[BEns]}] N.~Bourbaki，
  \emph{集合论〔Th\'eorie des ensembles〕}，第 I--IV 章，Hermann，1970。

  \item[\textup{[BLie]}] N.~Bourbaki，
  \emph{Lie 群与 Lie 代数〔Groupes et alg\`ebres de Lie〕}，第 II--III 章，Hermann，1970。

  \item[\textup{[Ca62]}] P.~Cartier，
  \emph{代数群与形式群〔Groupes alg\'ebriques et groupes formels〕}，第~87--111~页，载于：\emph{代数群论研讨会〔Colloque sur la th\'eorie des groupes alg\'ebriques〕}（Bruxelles，1962），Gauthier-Villars，1962。

  \item[\textup{[DG70]}] M.~Demazure，P.~Gabriel，
  \emph{代数群〔Groupes alg\'ebriques〕}，Masson \& North-Holland，1970。

  \item[\textup{[De72]}] M.~Demazure，
  \emph{p-可除群讲义〔Lectures on \(p\)-divisible groups〕}，
  Lect.\ Notes Math.\ \textbf{302}，Springer-Verlag，1972。

  \item[\textup{[Di55]}] J.~Dieudonn\'e，
  \emph{正特征域上的 Lie 群与 Lie 超代数〔Groupes de Lie et hyperalg\`ebres de Lie sur un corps de caract\'eristique \(p>0\)〕}，Comm.\ Math.\ Helv.\ \textbf{28} (1955)，87--118。

  \item[\textup{[Di73]}] J.~Dieudonn\'e，
  \emph{形式群论引论〔Introduction to the theory of formal groups〕}，Marcel Dekker，1973。

  \item[\textup{[Fo77]}] J.-M.~Fontaine，
  \emph{局部域上的 p-可除群〔Groupes \(p\)-divisibles sur les corps locaux〕}，Ast\'erisque \textbf{47--48}，Soc.\ Math.\ France，1977。

  \item[\textup{[Gr57]}] A.~Grothendieck，
  \emph{关于同调代数的若干问题〔Sur quelques points d'alg\`ebre homologique〕}，T\^ohoku Math.\ J.\ \textbf{9} (1957)，119--221。

  \item[\textup{[Gr74]}] A.~Grothendieck，
  \emph{Barsotti-Tate 群与 Dieudonn\'e 晶体〔Groupes de Barsotti-Tate et cristaux de Dieudonn\'e〕}，Presses Univ.\ Montr\'eal，1974。

  \item[\textup{[HS69]}] R.~G.~Heyneman，M.~E.~Sweedler，
  \emph{仿射 Hopf 代数 I〔Affine Hopf Algebras I〕}，J.\ Algebra \textbf{13} (1969)，194--241。

  \item[\textup{[Il85]}] L.~Illusie，
  \emph{依 A.~Grothendieck 的 Barsotti-Tate 群的形变〔D\'eformations de groupes de Barsotti-Tate, d'apr\`es A.~Grothendieck〕}，第~151--198~页，载于：\emph{算术线系讨论班：Mordell 猜想〔S\'eminaire sur les pinceaux arithm\'etiques: la conjecture de Mordell〕}，Ast\'erisque \textbf{127}，Soc.\ Math.\ France，1985。

  \item[\textup{[Ja65]}] I.~M.~James，
  \textup{[MM65]} 在 Math.\ Reviews 中的书评，MR0174052。

  \item[\textup{[La75]}] M.~Lazard，
  \emph{交换形式群〔Commutative formal groups〕}，
  Lect.\ Notes Math.\ \textbf{443}，Springer-Verlag，1975。

  \item[\textup{[LT65]}] J.~Lubin，J.~Tate，
  \emph{局部域中的形式复乘〔Formal complex multiplication in local fields〕}，Ann.\ of Math.\ \textbf{81} (1965)，380--387。

  \item[\textup{[LT66]}] J.~Lubin，J.~Tate，
  \emph{单参数形式 Lie 群的形式模空间〔Formal moduli for one-parameter formal Lie groups〕}，Bull.\ Soc.\ Math.\ France \textbf{94} (1966)，49--60。

  \item[\textup{[Ma91]}] A.~Masuoka，
  \emph{具有余交换余根基的 Hopf 代数〔On Hopf algebras with cocommutative coradicals〕}，J.\ Algebra \textbf{144} (1991)，451--466。

  \item[\textup{[Me72]}] W.~Messing，
  \emph{与 Barsotti-Tate 群相联系的晶体及其对阿贝尔概形的应用〔The crystals associated with Barsotti-Tate Groups: with applications to abelian schemes〕}，Lect.\ Notes Math.\ \textbf{264}，Springer-Verlag，1972。

  \item[\textup{[MM65]}] J.~W.~Milnor，J.~C.~Moore，
  \emph{论 Hopf 代数的结构〔On the structure of Hopf algebras〕}，Ann.\ of Math.\ \textbf{81} (1965)，211--264。

  \item[\textup{[Mi65]}] B.~Mitchell，
  \emph{范畴论〔Theory of categories〕}，Academic Press，1965。

  \item[\textup{[Ne75]}] K.~Newman，
  \emph{余交换 Hopf 代数中双理想与子 Hopf 代数之间的对应〔A correspondence between bi-ideals and sub-Hopf algebras in cocommutative Hopf algebras〕}，J.\ Algebra \textbf{36} (1975)，1--15。

  \item[\textup{[Po73]}] N.~Popescu，
  \emph{阿贝尔范畴及其对环与模的应用〔Abelian categories with applications to rings and modules〕}，Academic Press，1973。

  \item[\textup{[Sch90]}] H.~J.~Schneider，
  \emph{任意 Hopf 代数的主齐性空间〔Principal homogeneous spaces for arbitrary Hopf algebras〕}，Israel J.\ Math.\ \textbf{72} (1990)，第~1--2~期，167--195。

  \item[\textup{[Sw67]}] M.~Sweedler，
  \emph{只有一个群状元素的 Hopf 代数〔Hopf algebras with one grouplike element〕}，Trans.\ Amer.\ Math.\ Soc.\ \textbf{127} (1967)，第~3~期，515--526。

  \item[\textup{[Sw69]}] M.~Sweedler，
  \emph{Hopf 代数〔Hopf algebras〕}，Benjamin，1969。

  \item[\textup{[Tak72]}] M.~Takeuchi，
  \emph{Hopf 理想与子 Hopf 代数之间的对应〔A correspondence between Hopf ideals and sub-Hopf algebras〕}，Manuscripta Math.\ \textbf{7} (1972)，251--270。

  \item[\textup{[Tak79]}] M.~Takeuchi，
  \emph{相对 Hopf 模——等价性与自由性判据〔Relative Hopf modules--Equivalence and freeness criteria〕}，J.\ Algebra \textbf{60} (1979)，452--471。

  \item[\textup{[Ta67]}] J.~Tate，
  \emph{p-可除群〔\(p\)-divisible groups〕}，第~158--183~页，载于：\emph{局部域会议文集〔Proc.\ Conf.\ Local Fields〕}（Driebergen，1966）（T.~A.~Springer 编），Springer-Verlag，1967。

  \item[\textup{[We95]}] C.~A.~Weibel，
  \emph{同调代数引论〔An introduction to homological algebra〕}，Cambridge Univ.\ Press，1995。
\end{description}

% 第 VII_B 篇及第 VII 篇结束。 下一处不在本范围内的游标是第 VIII 篇分篇第 1 页、合并阅读本第 617 页。第一卷索引第 613--616 页不属于第 VII 篇。
